\section{Two ports on a polygon}\label{L:part8}

Glue both positive vacua at a slice with columns in cyclic order \([a,\mathcal X,b,\mathcal Y]\), attributing weight only to a \emph{single} loop containing both named ports \(a,b\) (at the interfaces in time of consecutive rows). Denote this sum by \(C_{ab}\). The identity below expresses this observable in the divided-slice colors.
Using \(m,n,\sigma,D,P,w\) for the regular variables \((\mathcal X,\mathcal Y)\) exactly as in \eqref{L:eq2} of the divided-slice calculation (\(q=e^{\pi i/4},\ \phi_+=\phi_-=\pi/4\)), we have
\begin{unnumberedproposition}[Two-port identity]
\phantomsection\label{L:prop:two-port-identity}
\begin{equation}
 Q(a,\mathcal X,b,\mathcal Y)^2 C_{ab} =
 \rho a b \prod_{i<j}D_{\sigma_i\sigma_j}
 \sum_{\sum \sigma_i t_{c_i}=0}
 \prod_{i<j}P^{\sigma_i\sigma_j}_{c_i c_j}(z_j/z_i)
 \prod_i w_{c_i}^{\sigma_i} f_{c_i}(a,z_i) f_{c_i}(b,z_i)\label{L:eq7}
\end{equation}
where \(f_c(p,z)=(p-q^r z)(p-q^{r'}z)\) for \((r,r')=(5,6),(3,5),(2,3)\) respectively.
\end{unnumberedproposition}

\begin{proof}
We compare the two sides after clearing the denominator $Q^2$.
The usual pair reductions and endpoint values determine most of the
polynomial. A slide across either mark controls the remaining cases,
where the two intervals differ in size by at most two. Two paired
specializations finish the balanced case.

\paragraph{Reduction and slide rules.} The denominator, fusions, removals and symmetry within each regular group on the left are just as before: the mark insertions do not interfere with those surgeries. Closed cycles must be killed except one passing both marks, which cannot be closed off during such vacuum contractions. Each mark variable contributes a factor of itself to the polynomial and saves one from its degree bound, by mandatory occupation in both vacua.

Both sides fulfill the same rules, including homogeneity \((m+n+2)(m+n+1)\), mark bounds, degrees \(\le2(m+n+1)\) per regular variable. Indeed the pole and endpoint checks of \eqref{L:eq2} carry through: for the three fusions \(+0,0-,+-\) (at high label position \(q^2x\), in the order low then high) the added pair factors \(f\) per mark \(p\) relative to the reduced one equal \((p+qx)^2\); for gap three the disappearing \(+-\) contributes \(d(p,x)^2\). At gap one the extra factors agree in the canceling terms. At the two endpoints the only label contributing is 0. There is an additional slide rule on both sides: when \(a=q^{\pm3} z\) with \(z\in\mathcal X\), one can instead take marks \(z,b\) and intervals \(\mathcal X-z,\mathcal Y+a\), and likewise with roles interchanged. Indeed arrange \(z\) adjacent to its mark and one of the vacua has the cap there, so incidence is equivalent. In \eqref{L:eq7} at \(a=q^3z\) only \(c(z)=+\) survives, exchanged with \(c(a)=-\); neutrality is preserved. The identities
\[
\begin{gathered}
 D_\varepsilon(z,u)P^\varepsilon_{+c}(u/z) f_c(a,u)
 =D_{-\varepsilon}(a,u)P^{-\varepsilon}_{-c}(u/a) f_c(z,u),\\
 a w_+ f_+(a,z) f_+(b,z)= z w_- f_-(z,a) f_-(b,a).
\end{gathered}
\]
by substitution verify equality; conjugation gives the reversed check (or exchange the two sides).



Start induction in total size. With no regular variables \eqref{L:eq7} is the two-site positive-vacuum component product \(\rho ab/(a+b)^2\). For one regular \(x\) in order \(a,x,b\), only the \(ab\)-cup contributes in each vacuum. In the lower vacuum its polynomial after clearing \(Q\sqrt{ab}^{\,-1}\) (multiplying by \(Q\) and dividing by \(\sqrt{ab}\)) is a constant times \((a-q^3x)(x-q^3 b)\) by adjacent caps and degree two, and the reflected one has reversed zeros. Thus the product after multiplying by \(Q^2\) has form const. times \(ab h(a,x)h(b,x)\), and extreme removal gives \eqref{L:eq7}.

\paragraph{The residual polynomial.}
Assume the identity at smaller total sizes. Subtract the two polynomial
numerators. The regular-variable reductions give zeros at the within-type
pair nodes, and matching removal values give a zero at each variable's
origin and cancel its leading coefficient. Their difference is therefore
\[
 ab\prod_i z_i\prod_{i<j,\ \sigma_i=\sigma_j} K(z_i,z_j)\,H(a,b;\mathcal X,\mathcal Y)
\]
with \(H\) polynomial of total degree \(4(m+n)-(m-n)^2\) and bounds \(4\mp2(m-n)\) in regular variables of either type. Thus only \(|m-n|\le2\) remain. The slide says at \(z=q^{\ell}a,\ \ell=\pm3\), \(z\in\mathcal X\), that
\[
 \prod_{x\in\mathcal X-z}K(z,x)\,H=\prod_{y\in\mathcal Y}K(a,y)\,
                    H(z,b;\mathcal X-z,\mathcal Y+a).
\]
All analogous variants hold.

\paragraph{Odd regular size.}
At odd regular size, say \(m=n+1,\ n\ge1\), \(H\) contains \(\prod_{y}h(a,y)h(b,y)\) by sliding a minority variable to imbalance three. Write \(U\) for the quotient with degrees at most two in all regular variables; similarly across the slide into opposite imbalance. Sliding a majority \(z=q^{\pm3}a\) now shows by divisibility in the reduced slide equation that \(U\) there is a scalar (in the mark variables) times \(\prod_y(a^2+y^2)\prod_{x\ne z}h(b,x)\); likewise at the two \(b\)-nodes. These two spectator products are linearly independent for generic
marks. Regard $U$ as a quadratic polynomial in $z$ with coefficients in
the vector space of spectator polynomials. Its four node values lie in
the span of those two products, so its coefficients do as well: any
linear functional vanishing on that span gives a quadratic vanishing at
four distinct points. Project now onto the first product modulo the
second. The resulting scalar quadratic vanishes at the two $b$-nodes,
so it is proportional to $h(b,z)$. The other projection is proportional
to $h(a,z)$. Restoring the factor in $z$ gives the two full products. This gives
\[
 U=A(a,b)\prod_y(a^2+y^2)\prod_x h(b,x)+B(a,b)\prod_y(b^2+y^2)\prod_x h(a,x).
\]
Both coefficients are homogeneous rational of degree 1 and regular wherever the products stay independent (in particular \(a\ne b\), including a single zero). Let \(A',B'\) be the analogous coefficients for the opposite imbalance. The slide equation after substitution now reads \(A(a,b)/h(b,a)=A'(q^\ell a,b)/h(b,q^\ell a)\) at both shifts. Thus the right-hand function is invariant under first-argument rotation by \(i\), and can have no poles (every putative pole divisor can be rotated clear of equality and of the two zeros of \(h\)). Having degree \(-1\) it vanishes. Likewise the \(b\)-slide gives \(B=0\).

\paragraph{Even regular size.}
For even regular size, at \(m=n+2\) the polynomial is independent of \(\mathcal X\). Sliding a majority site as above shows it contains \(\prod_y K(a,y)K(b,y)\); the quotient by degrees is a homogeneous quartic \(A_+(a,b)\), likewise \(A_-\) at opposite imbalance. In balanced size \(m=n\ge1\), we thus know
\[
 H|_{z=q^\ell a}=A_-(q^\ell a,b)\prod_y K(a,y)\prod_{x\ne z}K(b,x)
\]
for majority-to-be-removed \(z\in\mathcal X\). The simultaneous slide at \(y=q^{\ell'}b\) (separately comparing its formula for \(H\)) gives \(A_-(q^\ell a,b)/K(q^\ell a,b)=A_+(a,q^{\ell'}b)/K(a,q^{\ell'}b)\) for both shifts. Again invariance under quarter-rotation removes all poles, leaving one constant \(c\). Subtract the now known slide values by explicit extension. Zeros and degrees force
\[
 H=c\left(\prod_y K(a,y)\prod_x K(b,x)+\prod_y K(b,y)\prod_x K(a,x)\right)
     +d'\prod_i h(a,z_i)h(b,z_i).
\]
\paragraph{The two paired specializations.}
Pair all of \([a,\mathcal X]\) as high rapidities at positive gap three against \([b,\mathcal Y]\) in reversed order. Both vacua are simultaneous caps (peel at opposite interfaces), so \(a,b\) on distinct matched pairs cannot be on one loop. The sum in \eqref{L:eq7} vanishes as well: \(t(x)\le t(y)\) on each regular pair by cross zeros, while the regular partners of \(b,a\) have labels \(-,+\), contradicting neutrality. Thus \(c=0\).

Instead now take \(a=q^3b\) and pair \(x=q^3y\) for all regulars in reversed order. In \eqref{L:eq7} all surviving paired labels agree by neutrality. Distinct regular pair interactions decouple as at the end of the proof of \eqref{L:eq2}, with \(Q(a,b,\{z_i\})=(a+b)\prod d(z_i,b)Q(\{z_i\})\). Thus the predicted \(C_{ab}\) is the product over pairs of
\[
 \sum_c q^{2t_c}\,
 \frac{ f_c(a,x) f_c(b,x) f_c(a,y) f_c(b,y)}
 {d(x,b)^2d(y,b)^2}
 =
 v(s)v(1/s)+w(s)W(1/s)+W(s)w(1/s),\qquad s=a/x .
\]
Here lower-case and upper-case functions on the right are cell weights of the periodic calculation. The scalar equality follows by substituting \(a=q^3 b,\ x=a/s,\ y=b/s\) and using \(2i\sin\eta=s^{1/2}-s^{-1/2}\) in the weights; explicitly with \(d_s=(s-q^5)(s-q^6)\) we have \(d_s(v,w,W)=((s-q^3)(s-1),q^3(s-q^3)(s-q^2),q^5(s-q)(s-1))\), and the three terms for labels \(+,0,-\) give \(W(s)w(1/s), v(s)v(1/s), w(s)W(1/s)\). To check the diagram sum itself start with both vacua at simultaneous caps in order \(\mathcal X,a,b,\mathcal Y\), then exchange \(a\) back across each \(x\), with swaps \(B(a/x),B(x/a)\) in lower and reflected upper states. Gluing across \(\mathcal Y\) connects the original lower and upper ports at each \(x\) if occupied, by its two caps through the common port at \(y\); the exchanged \(x\)'s are glued directly. Original \(a\)'s connect via the occupied caps through \(b\). Both exchanged \(a\) lines (one in each chain of swaps) must be occupied all the way across the block of \(x\)'s with through connections at each step, to have a single loop through both initial and final marked connections. At each pair of swaps in that block, either both are single opposite-side arcs (\(v\)), or both double and using opposite pairings (\(w,W\) in either order); all such choices do give allowed diagrams. This proves exactly the product just found. At generic such specialization the factor multiplying \(d'\) is nonzero, so \(d'=0\), completing induction.
\end{proof}
